\section{Introduction}
\label{sec:introduction}

A noncommutative rational formula is a rooted tree whose leaves are variables or rational constants, whose binary gates are addition or ordered multiplication, and whose unary gates are inversion. Its size is its number of vertices, including leaves and inverse gates. At a tuple of square rational matrices, a scalar is interpreted as a scalar multiple of the identity. Evaluation is defined only when every inverse gate encounters an invertible matrix.

Call a formula \emph{admissible} if it has a defined rational matrix evaluation of some positive dimension, and \emph{nonzero} if one of its defined rational matrix evaluations is nonzero. These operational definitions suffice for the hitting-list theorem; its proof does not need a construction of the free skew field. In particular, algebraic cancellations do not enlarge a formula's domain of evaluation.

Identity testing asks whether an admissible formula is zero at every
defined evaluation. Scalar substitutions cannot answer this question:
for example, $x_1x_2-x_2x_1$ vanishes on all scalar tuples but is
nonzero on matrices. A matrix hitting list is a fixed finite list of
tuples that supplies a defined nonzero evaluation for each nonzero
formula in the prescribed class. Such a list permits testing through
evaluations without access to the formula tree. With inverse gates,
the same tuple must satisfy every domain condition imposed by that
tree, including conditions in branches that later cancel.
In oracle terminology, an undefined evaluation is reported as such;
only defined outputs are tested for zero.

\begin{theorem}\label{thm:main}
There are a deterministic Turing machine $G$ and fixed positive integers $C,k$ with the following property. Given $1^n,1^s$, where $n,s\geq1$, the machine outputs a finite list $\mathcal H_{n,s}$ of $n$-tuples of rational square matrices, in at most $C(n+s+1)^k$ bit operations. The entire output, with signed binary numerators and positive binary denominators, has at most $C(n+s+1)^k$ bits. For every nonzero admissible rational formula $\Phi$ in $n$ variables of size at most $s$, some tuple in $\mathcal H_{n,s}$ is a defined evaluation of $\Phi$ with invertible value.

The same list works for all such formulas, regardless of the bit lengths of their rational constants or their inverse depth. The generator receives no formula, oracle, randomness, or advice.
\end{theorem}

The invertible-value conclusion is stronger than the nonzero-value requirement. All tuples produced below even have one common dimension depending only on $n,s$.

\subsection{Context}
The algebraic setting comes from the theory of rational identities
and universal fields of fractions~\citep{Amitsur66,Cohn06}.
Affine-matrix representations of free-field elements appear in
Cohn and Reutenauer~\citep[Section~3.1]{CR94};
Kaliuzhnyi-Verbovetskyi and Vinnikov give a systematic account of
matrix evaluation domains and rational functions~\citep[Section~2]{KVV10}.
The distinction between a rational function and an expression
representing it matters here: equivalent expressions can impose
different conditions on their intermediate inverses. Our operational
formulation retains the original expression throughout.

For division-free computation, Nisan's coefficient-matrix
characterization made rank a central tool in noncommutative
branching-program complexity~\citep{Nisan91}.
Raz and Shpilka gave deterministic polynomial-time identity testing
with access to a formula or branching program~\citep{RS05}.
Bogdanov and Wee developed randomized black-box matrix evaluation
for bounded-degree noncommutative polynomials~\citep{BW05}, and
Forbes and Shpilka constructed deterministic quasipolynomial-size
matrix hitting sets for noncommutative algebraic branching programs,
a model including division-free formulas~\citep{FS13}.

For computation with division, Hrube\v{s} and Wigderson developed
the rational-formula model and its reduction to linear-matrix
singularity testing~\citep{HW15}. Their inverse-matrix realizations
and tests for the domains of gate subformulas are direct predecessors
of our certificate construction. Deterministic polynomial-time
algorithms with access to the input formula followed from operator
scaling and constructive noncommutative-rank
computation~\citep{GGOW16,GGOW20,IQS18}. These algorithms use the
given coefficients, with their encoding lengths included in the
bit model. Derksen and Makam established polynomial bounds on the
dimensions needed for an admissible rational formula to have a
defined evaluation, and derived randomized polynomial-time rational
identity testing~\citep[Section~6]{DM17}. Small witnesses and
algorithms that depend on the input coefficients leave open the
construction of a universal deterministic list from size bounds alone.

Arvind, Chatterjee, and Mukhopadhyay obtained deterministic
quasipolynomial-size hitting sets for rational formulas of inversion
height two~\citep[Theorem~1.1]{ACM22}.
Their subsequent general result handles unrestricted inverse
nesting with quasipolynomial-size rational matrix hitting
sets~\citep{ACM26}; the quantitative bounds used in this comparison
are Theorem~2, Corollary~3, and the paragraph following Remark~52
of the specified preprint version~\citep{ACM25}. Thus rational matrix outputs and arbitrary
inverse nesting already occur in the preceding literature.
Invertible evaluations are also used in the division-algebra methods
of these works~\citep{ACM22,ACM25}. Their general treatment states a
depth reduction that preserves the input formula's evaluation
domain~\citep[Fact~1]{ACM25}. Theorem~\ref{thm:main} gives polynomial
bounds on construction time, total binary output length, and matrix
dimension.
Its complexity bound concerns generation of the list; the cost of
evaluating an unknown formula with arbitrarily large rational
constants is not included.

The result resolves, over $\Q$, the explicit polynomial-size hitting-set
problem for this formula model, and gives the polynomial matrix
dimension asked for in the conclusion of~\citet{ACM25}.
The proof of Theorem~\ref{thm:main} below establishes all the algebraic and algorithmic
ingredients it uses.

The proof also draws on established methods. The multiplicity
argument for explicit subspace designs of Guruswami and
Kopparty~\citep{GK16}, and the connection between subspace designs,
rank condensers, and dimension expansion studied by Forbes and
Guruswami~\citep{FG15}, provide context for our
construction. The underlying Wronskian calculation is
classical~\citep{BD10}. The common-subspace obstruction to
invertibility belongs to noncommutative-rank
theory~\citep{FR04}, while cyclic division algebras and rank
divisibility underlie the regularity methods of Ivanyos, Qiao,
and Subrahmanyam~\citep{IQS17,IQS18}. We prove the precise expansion,
deformation, and domain statements needed here, without importing
the stronger general algorithms from those works.

\subsection{The construction and its logic}

Write $\Mat_q(F)$ for the $q$ by $q$ matrices over a field $F$.
The bridge from formula evaluation to linear algebra is an
\emph{affine linear pencil}, a matrix
\[
 L(x)=A_0+\sum_{i=1}^n A_ix_i,\qquad A_i\in\Mat_q(\Q),
\]
of order $q$. Its evaluation at matrices $X_i\in\Mat_e(E)$ over an
extension field $E/\Q$ is $A_0\otimes I_e+\sum_iA_i\otimes X_i$.
Call the pencil \emph{feasible} if one such evaluation is invertible.
Tensor products and ranks are taken over the indicated field.

A size-$s$ formula gives at most $s+1$ pencils of order at most
$B=2s+1$. Their simultaneous invertibility certifies that every
inverse occurrence in the original tree is defined and that the
final value is invertible. A nonzero admissible formula has a
matrix evaluation over an extension field satisfying all these
conditions. Section~\ref{sec:encoding} proves this reduction, including
the exact conditions at the inverse gates. The construction therefore
seeks a universal rational list on which each such finite family of
feasible pencils is simultaneously invertible.

Two obstacles remain: obtaining full rank at matrices independent of
the pencil coefficients, and keeping all conditions valid when the
parameters become rational numbers. The proof handles them in the
following order.
\begin{enumerate}
\item Section~\ref{sec:expansion} constructs $n+1$ rectangular
matrices over $\Q[t]$. They act on coefficient vectors of polynomials,
reversing the Taylor expansion about the distinct points $0,\ldots,n$
and weighting the Taylor coefficients by powers of $t$.
A Wronskian bound shows that, for every nonzero subspace of dimension
at most $B$, the sum of its images has dimension greater than $n$
times its dimension, even when the subspace has Laurent-series
coefficients. The resulting test for a feasible
pencil is injective. Padding the test matrices to a common square
order $M$ gives matrices $R_i(t)$ for which every feasible pencil of
order $1\leq q\leq B$ satisfies
\[
 \rank\Bigl(\sum_{i=0}^n A_i\otimes R_i(t)\Bigr)>(q-1)M.
\]
\item Section~\ref{sec:division} uses two further parameters $p,r$
to lift these matrices into a division algebra represented by
$M\times M$ matrices. Ranks of block matrices over this algebra are
multiples of $M$, so the strict inequality forces full rank.
Dividing on the left by the lifted $R_0$ restores the constant term
of an affine pencil. A separate application of the lift turns any
defined nonzero formula evaluation into a defined invertible one
over an extension field. This supplies the feasibility assertion
needed for the certificates in Section~\ref{sec:encoding}.
\item Section~\ref{sec:generator} represents multiplication in a
finite-dimensional rational algebra by rational matrices, thereby
replacing the root of unity used in the lift. All root components
must be invertible simultaneously. For any fixed formula, the
product of their determinant conditions is a nonzero polynomial
of bounded degree in $t,p,r$. A fixed three-dimensional integer
grid detects it. Exact rational arithmetic bounds both the
generation time and the total binary output length.
\end{enumerate}

The witnesses, subspaces, and formula-specific determinant products in this proof certify the construction. None is computed by $G$. In particular, the dimension of an initial witness for a formula has no effect on the generator's resource bounds.

The construction also hits every feasible rational affine pencil
of bounded order; see Corollary~\ref{cor:pencil-hitting}.
In particular, it answers the rational-coefficient version of the
black-box \textsc{Singular} construction question for homogeneous
pencils posed by Garg, Gurvits, Oliveira, and
Wigderson~\citep[Section~6, p.~280]{GGOW20}.
Section~\ref{sec:rank} then treats rectangular affine pencils.
Their \emph{noncommutative rank} is their rank over the free skew
field. Corollary~\ref{cor:rank-detecting-list} constructs a universal
nonadaptive list whose normalized evaluation ranks recover this rank
for every rectangular rational affine pencil with prescribed row,
column, and variable counts. It permits exact matrix-value or
ordinary-rank responses and accounts for their processing costs.
